\documentclass[12pt]{article}
\usepackage[T1]{fontenc}
\usepackage[utf8]{inputenc}
\usepackage[brazil]{babel}
\usepackage[margin=1in]{geometry}
\usepackage{newtxtext,newtxmath}
\usepackage{ragged2e}
\usepackage[normalem]{ulem}
\setlength{\parindent}{0pt}
\linespread{1.15}
\setlength{\parskip}{0.7\baselineskip}
\emergencystretch=3em
\sloppy
\begin{document}
\justifying
\noindent \textbf{Processo n.: }\uline{5056703-85.2023.4.02.5101}

{\Large\bfseries \textbf{PU-TEMA TNU-SOBRESTAMENTO}}\\

\noindent Verifica-se que a questão objeto da controvérsia jurídica suscitada pelo recurso manejado pela parte recorrente encontra-se afetada ao Tema n. 172 da jurisprudência da TNU – PU no processo n. PEDILEF 0514224-28.2017.4.05.8013/AL, pendente de julgamento (afetado em 29/05/2018 26/10/2018 - manutenção da afetação) –, por meio do qual se elucidará a seguinte questão:

\noindent \textit{"Saber se é possível ou não aplicação da regra prevista no art. 29, I e II, da Lei 8.213/91, quando mais favorável que a regra de transição prevista no art. 3º da Lei 9.876/99."}

\noindent Nesse contexto, determino o \textbf{SOBRESTAMENTO} do PU interposto nestes autos – o qual versa sobre a mesma matéria –, na forma do\textbf{\uline{ art. 14, II, b}}, da Resolução CJF n. 586/2019 c/c 1.030, III, do CPC, até o julgamento final do representativo da controvérsia, o PU no processo n. PEDILEF 0514224-28.2017.4.05.8013/AL (\textbf{Tema n. }\textbf{172 }\textbf{TNU}).

\noindent Intimem-se. Cumpra-se.
\end{document}
